\section{Introduction} \label{sec:intro}
Despite rapid progress in speech language models, such as GPT-4o Voice Mode~\cite{hurst2024gpt} and Gemini Live~\cite{comanici2025gemini}, the naturalness of synthesized speech still lags behind human audio, with noticeable deficits in human-likeness, prosody, and emotional nuance, thus limiting the level of human-like engagement in real-worl applications. Bridging this gap requires a reliable evaluator of speech naturalness to guide the improvement of speech generation systems.

Traditional evaluators, including MOS predictors~\citep{lo2019mosnet, saeki2022utmos, baba2024utmosv2, wang2024uncertainty, samos2024, hoq2025fusemos}, enable scalable speech evaluation but ultimately reduce perceptual judgments to a single scalar score. While effective for large-scale benchmarking, such metrics provide limited insight into \emph{why} a given speech sample is perceived as unnatural. Speech naturalness is inherently multi-dimensional, including factors such as expressiveness, intonation, and pacing. Forcing a single regressor to capture all these aspects often leads to poor generalization~\citep{lo2019mosnet, saeki2022utmos}. As a result, scalar predictors can yield unreliable reward signals when used for online reinforcement learning. In the broader text-based LLM literature, \emph{Generative Reward Models} (GRMs) have recently emerged as a promising alternative to scalar reward predictors. Rather than producing only a numerical score, GRMs generate explicit reasoning traces alongside their evaluations, more closely aligning with how humans justify judgments~\citep{mahan2024generative, liu2025inference, zhang2024generativeverifiers, liang2025grmcriteria}.


However, generative reward modeling remains relatively under-explored in the speech domain. AudioJudge~\cite{manakul2025audiojudge} evaluates the ability of frontier speech LLMs to judge audio quality and paralinguistic attributes through prompt engineering, revealing that most models struggle to produce reliable judgments. WavReward~\cite{ji2025wavreward} finetunes a GRM to assess specific speech attributes such as age, gender, pitch, and volume. Concurrent work likeSpeechJudge~\cite{zhang2025speechjudge} extends generative reward modeling to judge speech naturalness. While both SpeechJudge and our approach aim to construct synthetic CoT reasoning, they differ fundamentally in how the CoT traces are produced. SpeechJudge synthesizes CoT using a teacher speech LLM directly conditioned on the raw audio signal, resulting in reasoning that is mediated by implicit acoustic representations the speech LLM can capture. In contrast, we explicitly extract interpretable paralinguistic acoustic cues via signal processing and generate CoT using a text LLM conditioned on these features. This explicit grounding enables systematic ablation of which acoustic factors influence the predicted naturalness.  Moreover, although GRMs have been used to improve text-to-speech (TTS) systems~\cite{zhang2025speechjudge}, no prior work has leveraged these models to improve the naturalness of a full-duplex speech LLM via online reinforcement learning.

To this end, we introduce \textbf{GSRM}, a \emph{Generative Speech Reward Model} tailored for speech RLHF, with a focus on speech naturalness evaluation. GSRM explicitly decomposes speech judgment into two stages: (1) extracting interpretable paralinguistic acoustic features from raw audio, and (2) generating feature-grounded CoT reasoning traces to produce final evaluations. We summarize our main contributions as follows:
\begin{enumerate}[leftmargin=*]
    \item \textbf{High-quality Human Feedback Data.} In Section~\ref{sec:human_data}, we curate a large-scale human feedback dataset for speech naturalness, comprising 6.5K unique audio samples generated by controllable TTS systems. In addition, we collect 490 dialogue samples generated by an internal speech LLM to emulate real-world interactive scenarios.

    \item \textbf{Acoustic Feature–grounded Reasoning.} Unlike existing approaches, GSRM explicitly extracts paralinguistic acoustic features from speech and reasons over this evidence to judge audio quality. Section~\ref{sec:gsrm} introduces a scalable CoT synthesis pipeline that enables the generation of high-quality, feature-grounded reasoning trajectories.

    \item \textbf{Strong Empirical Performance.} 
    GSRM achieves human-level performance and
    substantially outperforming exisiting methods on naturalness prediction. Moreover, we take a first step toward online RLHF for speech LLMs. In Section~\ref{sec:rlhf}, by leveraging GSRM as a verifier, we successfully improve the naturalness of a speech LLM through online reinforcement learning.
\end{enumerate}
